\documentclass{subfiles}
\begin{document}
    We now present some results that are at the core of the codes we present later.
        Though we try to be as formal as possible, we also try to make the proof 
        undestandable to those with little-to-no mathematical background.

    \begin{theorem}\label{Thm:6.1}
        Let \(\C*\) be linear code defined by a parity matrix \(H\) and
            let \(d\) its minimal distance. Then, 
            \begin{enumerate}
                \item Every \(d\) columns of \(H\) are linearly dependent, and 
                \item Every \(d - 1\) columns of \(H\) are linear independent.
            \end{enumerate}
    \end{theorem}
    \begin{proof*}
        \begin{enumerate}
            \item Recall that the minimal distance of a code, is also the minimal weight
                of distinct words in it. It follows that 
                \[
                    \exists \bv{w} \in \C* \text{ s.t. } \bv{w} = 
                        (0, \ldots, 0, \alpha_{1}, 0, \ldots, 0, \ldots, 0, \alpha_{d}, 0, \ldots, 0)
                \]
                where \(\alpha_{1}, \ldots, \alpha_{d} \ne 0\). By hypotesis
                \(H \bv{w}^{t} = \bv{0}\), hence there exist linear combination 
                of columns of \(H\) that is the zero-vector. Thus, we have \(d\) linearly 
                dependent columns.

            \item Trivial. By contradiction, if \(d - 1\) columns would be linearly dependent 
                then by 1., there exists a word \(\bv{w} \in \C*\) with \(d - 1\) non 
                zero components such that \(H \bv{w}^{t} = \bv{0}\). 
                But then the minimal weight of \(\bv{w}\) is \(d - 1\) and not \(d\)
                as supposed.\qed
        \end{enumerate}
    \end{proof*}

    \begin{theorem}
        Let \(\C*\) be a linear code and let \(d\) be its minimal distance. Then,
            \begin{enumerate}
                \item The code can detect up to \(d\) errors, and 
                \item Up to \(\sfrac{(d - 1)}{2}\) errors can be corrected.
            \end{enumerate}
    \end{theorem}

    \begin{proof*}
        \begin{enumerate}
            \item Let \(\bv{w} \in \C*\) be the word transmitted,
                and let \(\bv{z}\) be the one recived. Recall that both,
                sender and reciver, know the code. Assume that during the transmission 
                \(e\) errors occured. If \(e < d\), then \(\bv{z}\)
                is not a valid word, in fact we would have \(d(\bv{w}, \bv{z}) = e < d\). 
                On the other hand, if \(e > d\), we can say nothing about \(\bv{z}\)
                and it would be accepted as valid.

            \item By contradiction, suppose that \(\bv{w}_{1} \ne \bv{w}_{2}\) are 
                two words with distance at most \(\sfrac{(d - 1)}{2}\) from \(\bv{z}\).
                That is, 
                \[
                    d(\bv{w}_{1}, \bv{z}) \le \frac{d - 1}{2}, 
                        \qquad d(\bv{w}_{2}, \bv{z}) \le \frac{d - 1}{2}.
                \]
                By the trinangular inequality, we have 
                \[
                    d(\bv{w}_{1}, \bv{w}_{2}) \le 
                        d(\bv{w}_{1}, \bv{z}) + d(\bv{w}_{2}, \bv{z})
                        = \frac{d - 1}{2} + \frac{d - 1}{2} = d - 1,
                \]
                but, by hypotesis, \(\bv{w}_{1}, \bv{w}_{2}\) have distance \(d\).
                Meaning that \(\bv{w}_{1} = \bv{w}_{2}\).\qed
        \end{enumerate}
    \end{proof*}

    \begin{theorem}[Singleton bound]
        Let \(\C*\) be a \([n, k, d]_{q}\) linear code defined by a parity matrix \(H\). 
            Then, it holds 
            \[
                d \le n - k + 1.
            \]
    \end{theorem}

    \begin{proof*}
        Recall that by the dimensional relation \(r = \rk{H} = n - k\). By 
            \emph{Theorem \ref{Thm:6.1}} \(d - 1\) columns of \(H\) are linearly independent.
            Combining the two, we have 
            \[
                d - 1 \le r \iff d - 1 \le n - k \iff d \le n - k + 1.
            \]\qed
    \end{proof*}

    A code that satisfies the Singleton bound is said to be optimal.

    \begin{theorem}[Hamming inequality]
        Let \(\C*\) be a \([n, k, d]_{q}\) linear code. Then, for any transmission errors
        \(e \le \sfrac{(d - 1)}{2}\), it holds
        \[
            q^{e} \left\lparen
                \sum_{i = 0}^{e}(q - i)^{i}\binom{n}{i}
            \right\rparen \le q^{n}.
        \]
    \end{theorem}

    \begin{proof*}
        Trivial. Let \(\bv{w}\) be the word transmitted, we have that:
        \begin{itemize}
            \item If \(e = 0\), we recived the correct word.
            \item If \(e = 1\), we have to consider the wordis with distance \(1\) from 
                \(\bv{w}\), which are exactly \((q - 1)n\).
            \item If \(e = m\), with \(m \le \sfrac{(d - 1)}{2}\), we have \((q - m)^{m}\binom{n}{m}\) 
                word at distance \(m\) from \(\bv{w}\).
        \end{itemize}
        Hence summing all up, we get
        \[
            q^{e} \left\lparen
                \sum_{i = 0}^{e}(q - i)^{i}\binom{n}{i}
            \right\rparen \le q^{n}.
        \]\qed
    \end{proof*}

    A code satisfying Hamming inequality is said to be perfect.
\end{document}
